\clearpage
\fancyhead[L]{\small PRÄZESSION.}
\SourceLeaf{566}
\begin{center}\small Göttingische gelehrte Anzeigen. Stück 19. Seite 188..189. 31. Januar 1814.\\[3pt]\rule{.72\linewidth}{.3pt}\end{center}

Die Königliche Akademie der Wissenschaften in Berlin hatte die möglichst genaue Bestimmung der Größe der Präzession zur Preisaufgabe gemacht. Sie hat Herrn \textsc{Bessel} für seine Arbeit darüber den Preis zuerkannt. Tatsächlich hätte niemand diese Bestimmung mit größerem Erfolg vornehmen können als Herr \textsc{Bessel}, der seit sechs Jahren an einer vielseitigen Auswertung der \textsc{Bradley}-schen Beobachtungen gearbeitet hat. Dazu hat er 4585 Sternbestimmungen verwendet. Wir glauben, den Astronomen einen wichtigen Dienst zu leisten, indem wir hier das von ihm gefundene Endergebnis für eine Größe mitteilen, die sie täglich benötigen.

\begin{center}\fontsize{8}{11}\selectfont\renewcommand{\arraystretch}{1.18}\begin{tabular}{@{}l@{\hspace{4pt}}l@{}}
Lunisolarpräzession $.\;.\;.\;.\;.\;.\;.\;.\;.\;.\;.\;.$ & $= 50''35330 - 0''0002435890\,(t-1800)$ \\
Beobachtete allgemeine Präzession $.\;.\;.\;.\;.\;.$ & $50.18924 + \mathrm{c}.0002442966\,(t-1800)$ \\
Konstante bei der Präzession in Rektaszension & $46.01058 + 0.0003590677\,(t-1800)$ \\
Konstante bei der Präzession in Deklination $.\;.$ & $20.04966 - 0.0002135621\,(t-1800)$ \\
\end{tabular}\end{center}

Die Bewegungen an einigen Punkten des Himmels sollen nach einer Beobachtung von \textsc{Triesnecker} auf eine ungleiche Präzession hindeuten. Herr \textsc{Bessel} hat keine deutliche Spur dieser Erscheinung bemerkt.
